\begin{statementbox}

\textbf{\textcolor{CStatement}{条件}}
\begin{description}[style=nextline, leftmargin=2.8em, labelsep=0.5em, itemsep=0.35em]
\item[\textbf{\textcolor{CStatement}{条件 1（整数与角）：}}] $k\in\mathbb{Z}$，$\alpha$ 为任意角。诱导公式本身对任意角成立；但使用口诀“奇变偶不变，符号看象限”判号时，通常把 $\alpha$ 暂看作锐角或参考角。
\item[\textbf{\textcolor{CStatement}{条件 2（定义域）：}}] 当公式涉及 $\tan$ 或 $\cot$ 时，必须保证等式左右两边均有意义：出现 $\tan\theta$ 要求 $\cos\theta\neq0$，出现 $\cot\theta$ 要求 $\sin\theta\neq0$。
\end{description}

\textbf{\textcolor{CStatement}{结论}}
\begin{description}[style=nextline, leftmargin=2.8em, labelsep=0.5em, itemsep=0.45em]
\item[\textbf{\textcolor{CStatement}{结论一（对称与周期公式）：}}]
\textbf{\textcolor{CStatement}{直观描述：}} 角 $-\alpha$, $\pi\pm\alpha$, $2k\pi+\alpha$ 可由单位圆的轴对称、中心对称与周期性得到，函数名通常不变，符号由对应象限或对称方式决定。
\[
\begin{aligned}
\sin(-\alpha)&=-\sin\alpha, & \cos(-\alpha)&=\cos\alpha,\\
\tan(-\alpha)&=-\tan\alpha, & \cot(-\alpha)&=-\cot\alpha;
\end{aligned}
\]
\[
\begin{aligned}
\sin(\pi-\alpha)&=\sin\alpha, & \cos(\pi-\alpha)&=-\cos\alpha,\\
\tan(\pi-\alpha)&=-\tan\alpha, & \cot(\pi-\alpha)&=-\cot\alpha;
\end{aligned}
\]
\[
\begin{aligned}
\sin(\pi+\alpha)&=-\sin\alpha, & \cos(\pi+\alpha)&=-\cos\alpha,\\
\tan(\pi+\alpha)&=\tan\alpha, & \cot(\pi+\alpha)&=\cot\alpha.
\end{aligned}
\]
一般地，$\sin(2k\pi+\alpha)=\sin\alpha$，$\cos(2k\pi+\alpha)=\cos\alpha$；并且 $\tan(k\pi+\alpha)=\tan\alpha$，$\cot(k\pi+\alpha)=\cot\alpha$，其中 $k\in\mathbb{Z}$，且两边均需有定义。

\item[\textbf{\textcolor{CStatement}{结论二（变名公式）：}}]
\textbf{\textcolor{CStatement}{直观描述：}} 角 $\frac{\pi}{2}\pm\alpha$ 会使正弦与余弦互换、正切与余切互换；这就是“奇变”。符号仍由整体角所在象限决定。
\[
\begin{aligned}
\sin\left(\frac{\pi}{2}-\alpha\right)&=\cos\alpha, &
\cos\left(\frac{\pi}{2}-\alpha\right)&=\sin\alpha,\\
\tan\left(\frac{\pi}{2}-\alpha\right)&=\cot\alpha, &
\cot\left(\frac{\pi}{2}-\alpha\right)&=\tan\alpha;
\end{aligned}
\]
\[
\begin{aligned}
\sin\left(\frac{\pi}{2}+\alpha\right)&=\cos\alpha, &
\cos\left(\frac{\pi}{2}+\alpha\right)&=-\sin\alpha,\\
\tan\left(\frac{\pi}{2}+\alpha\right)&=-\cot\alpha, &
\cot\left(\frac{\pi}{2}+\alpha\right)&=-\tan\alpha.
\end{aligned}
\]

\item[\textbf{\textcolor{CStatement}{口诀版（统一记法）：}}]
若角写成 $\frac{k\pi}{2}\pm\alpha$，$k\in\mathbb{Z}$，则 $k$ 为偶数时函数名不变，$k$ 为奇数时正余弦互换、正余切互换；右边符号按整体角所在象限确定。注意：口诀判号时的 $\alpha$ 是参考角，公式中的 $\alpha$ 仍可为任意角。
\end{description}

\medskip
\textbf{\textcolor{CStatement}{参数限制：}} 含 $\tan$、$\cot$ 的公式必须逐条检查定义域。不能只记形式而忽略分母不为零，否则在方程化简中可能引入增根或错误结论。
\end{statementbox}
